\clearpage
\chapter[LITERATURE REVIEW]{Literature Review}
\label{ch:lit}

\section{Introduction}
This project draws on five bodies of work: anomaly detection on time series, root-cause localisation in microservices, models of how failure spreads, evaluation practice, and the observability tools that supply the data. Each claim about a cited work restates what that work documents. Where we compare methods ourselves, we say so.

\section{Review of Existing Work}

\subsection{Anomaly detection}
Isolation Forest \cite{liu2008isolation,liu2012isolation} isolates an outlier with a few random cuts and needs no density estimate, which suits per-service streams of mixed, skewed metrics. It powered the first version of this engine, and Chapter~\ref{ch:results} shows where it breaks: a change of workload. Network monitoring met the same problem earlier. Traffic on many links follows a few shared causes, so researchers model the readings as a low-dimensional subspace and call anomalous whatever falls outside it \cite{lakhina2004diagnosing,jackson1979control}. Distance and density methods \cite{ramaswamy2000efficient,breunig2000lof}, the histogram score HBOS \cite{goldstein2012hbos} and Hotelling's multivariate control chart \cite{hotelling1931,mahalanobis1936} with a shrinkage covariance \cite{ledoit2004} are the obvious alternatives; Chapter~\ref{ch:results} tests all of them. The median and the median absolute deviation give robust location and scale \cite{hampel1974,rousseeuw1993}. They underpin our rule layer and, with a tail-aware spread \cite{rousseeuw1993alternatives}, the load-aware model. Surveys of streaming and learned detection \cite{gama2014survey,chandola2009survey,ahmad2017unsupervised} and deployed monitoring detectors \cite{laptev2015egads,su2019omnianomaly} describe our setting: unlabelled streams, drifting baselines and costly false alarms.

\subsection{Evaluation of anomaly detectors}
Several studies fault the common habit of scoring detectors point by point and crediting a whole event when any point inside it is flagged \cite{wu2021flawed,kim2022rigorous}. They argue for range-based precision and recall that respect an event's duration \cite{tatbul2018precision}. We follow them. Chapter~\ref{ch:method} scores at the event level: a fault counts as detected only when an alarm names a truly involved service while the fault is active, and a false alarm is an episode outside every fault window.

\subsection{Root-cause localisation in microservices}
MonitorRank \cite{kim2013monitorrank} introduced a personalised random walk over the call graph, biased by how anomalous each service looks, and ranked services by the walk's stationary share. MicroRCA \cite{wu2020microrca}, CloudRanger \cite{wang2018cloudranger}, Microscope \cite{lin2018microscope}, AutoMAP \cite{ma2020automap} and MicroRank \cite{yu2021microrank} extended it with anomaly-weighted transitions, causal-graph construction and trace spectra. Causal approaches recover structure with the PC algorithm or hierarchical discovery \cite{ikram2022rcd,li2022circa,xin2023causalrca}. Learned methods localise with graph autoencoders and fuse metrics, traces and logs \cite{sun2024deephunt,zhang2023diagfusion,lee2023eadro,wang2021groot}. Two surveys place these methods in context: one on AIOps failure management \cite{notaro2021survey}, one on anomaly and root-cause work for microservices \cite{soldani2022survey}. Most need traces, labelled failures or training, and some need accelerated hardware. The random-walk family needs only a call graph and an anomaly score, so this project builds on it.

\subsection{Cascading failure and change detection}
The independent-cascade model of influence spread \cite{kempe2003influence} underlies our propagation model, with one change: we estimate the edge probabilities online instead of assuming them. Proper scoring rules, reliability diagrams and expected calibration error judge probabilistic forecasts \cite{brier1950verification,niculescu2005predicting,guo2017calibration}. Classical change-point detection attacks detection latency \cite{page1954cusum,adams2007bocpd}. PageRank \cite{page1999pagerank} and reciprocal-rank fusion \cite{cormack2009rrf} inform the ranking, Granger causality \cite{granger1969} and transfer entropy \cite{schreiber2000te} inform the study of call direction, and random forests and CART \cite{breiman2001rf,breiman1984cart} support the fault-type classifier.

\subsection{Benchmarks and chaos engineering}
Open benchmark applications such as DeathStarBench \cite{gan2019deathstar} and Train-Ticket \cite{zhou2018trainticket} supply realistic systems to test on, and chaos engineering \cite{basiri2016chaos} supplies the method for breaking them on purpose. Our injector follows that practice with one addition: it keeps its ledger private, so nothing inside the system under test carries a label.

\subsection{Observability tooling}
Production observability combines three kinds of tool. \emph{Metric systems} scrape numbers from endpoints that each service exposes. \emph{Distributed tracing} follows a request across services, and the application or a proxy in front of it must propagate trace identifiers. \emph{Application performance monitoring} puts agents inside the runtime. Each gives an accurate picture once the instrumentation exists, and none works on a stack nobody will modify. Service meshes supply topology without code changes but need a sidecar in every pod. The site-reliability literature \cite{beyer2016sre} locates the hard part elsewhere: collecting signals is easy, deciding what to look at first is not. Root-cause ranking targets exactly that gap.

\begin{table}[!htbp]
\centering
\caption{Review of existing work (what each documents, what it needs, and where it stops).}
\label{tab:litreview}
\small
\renewcommand{\arraystretch}{1.25}
\begin{tabularx}{\textwidth}{>{\raggedright\arraybackslash}p{2.5cm}>{\raggedright\arraybackslash}X>{\raggedright\arraybackslash}p{3.1cm}>{\raggedright\arraybackslash}p{3.3cm}}
\toprule
\textbf{Work} & \textbf{Approach} & \textbf{Needs} & \textbf{Limitation for our setting}\\
\midrule
Isolation Forest \cite{liu2008isolation} & Isolate outliers with random trees; score by path length & Unlabelled history & Gives a score, not a cause; blind to the workload\\
Subspace residual \cite{lakhina2004diagnosing} & Model the shared causes of many signals as a subspace; score the residual & A matrix of signals & Learns the subspace from history; slow to adapt to a new load pattern\\
MonitorRank \cite{kim2013monitorrank} & Personalised random walk over the call graph & Call graph, anomaly score & Direction and edge weights are inputs\\
MicroRCA \cite{wu2020microrca}, CloudRanger \cite{wang2018cloudranger}, Microscope \cite{lin2018microscope}, AutoMAP \cite{ma2020automap} & Extend the walk with anomaly-weighted transitions and graph construction & A graph and monitoring data (sources differ by method) & The graph is built from data the host alone does not give\\
MicroRank \cite{yu2021microrank} & Trace spectrum analysis with PageRank & Request traces & Needs tracing in every service\\
RCD \cite{ikram2022rcd}, CIRCA \cite{li2022circa} & Causal discovery or structural causal graph for localisation & Metrics, structure & Assumes a faithful causal model\\
Eadro \cite{lee2023eadro}, DeepHunt \cite{sun2024deephunt}, DiagFusion \cite{zhang2023diagfusion} & Learned fusion of metrics, traces, logs; graph neural or autoencoder models & Labelled or curated failures; GPU & Not auditable; training data needed\\
Independent cascade \cite{kempe2003influence} & Probabilistic spread over a graph & Edge probabilities & Probabilities must be estimated\\
Change-point tests \cite{page1954cusum,adams2007bocpd} & Detect a shift in a stream & A stream & Alone, no cause or graph\\
\bottomrule
\end{tabularx}
\end{table}

\section{Comparative Analysis}
Table~\ref{tab:approaches} compares the approaches on the properties that decide whether a method can run on an unmodified Docker stack. The proposed system is not the most accurate method one could build. It switches on in minutes, explains itself, and reports its own failures.

\begin{table}[!htbp]
\centering
\caption{Approaches compared on practical properties (\checkmark\ yes, -- no, $\sim$ partly).}
\label{tab:approaches}
\small
\renewcommand{\arraystretch}{1.25}
\begin{tabularx}{\textwidth}{>{\raggedright\arraybackslash}p{4.2cm}>{\centering\arraybackslash}p{1.45cm}>{\centering\arraybackslash}p{1.45cm}>{\centering\arraybackslash}p{1.45cm}>{\centering\arraybackslash}p{1.45cm}>{\centering\arraybackslash}p{1.45cm}}
\toprule
\textbf{Property} & \textbf{Metric tools} & \textbf{Tracing} & \textbf{Learned fusion} & \textbf{Walk RCA} & \textbf{This project}\\
\midrule
No change to applications & -- & -- & -- & $\sim$ & \checkmark\\
Works without labelled failures & \checkmark & \checkmark & -- & \checkmark & \checkmark\\
Recovers the call graph itself & -- & \checkmark & $\sim$ & -- & \checkmark\\
Learns call direction & -- & \checkmark & -- & -- & \checkmark\\
Per-service learned baseline & $\sim$ & -- & \checkmark & $\sim$ & \checkmark\\
Ranks a root cause & -- & $\sim$ & \checkmark & \checkmark & \checkmark\\
Predicts what fails next & -- & -- & -- & -- & \checkmark\\
Needs a GPU & -- & -- & \checkmark & -- & --\\
Every score can be audited & \checkmark & \checkmark & -- & $\sim$ & \checkmark\\
\bottomrule
\end{tabularx}
\end{table}

\section{Research and Design Gap}
Four gaps motivate the design.
\begin{enumerate}
\item \textbf{Observation without instrumentation.} The graph that random-walk and causal methods assume is usually given; recovering it, with direction, from the host alone is rarely addressed.
\item \textbf{Facts are not degrees.} A stopped container is not ``very anomalous''; it is a fact. Methods that fold container state into one score lose the most reliable signal there is, and our ablation shows how much that costs.
\item \textbf{Propagation is assumed, not learned.} Cascade prediction normally uses fixed edge probabilities. Learning them from the system's own incidents, with a prior that makes rarely used edges behave sensibly, is cheap and measurable.
\item \textbf{Evaluation is optimistic.} Point-adjusted scoring, tuned thresholds and a single data set flatter results; we use event-level scoring, a held-out split and paired tests, and report what failed.
\end{enumerate}

\section{Summary}
Existing root-cause methods are accurate but assume instrumentation, labels or hardware that small teams do not have, and are often evaluated optimistically. The proposed system, described next, is instrumentation-free, unsupervised and auditable.
